\chapter[\Acl*{ArBIB}]{\Acl*{ArBIB}}
\label{sec:arbib-chapter}

While \ac{ArBIB} preparation is often used for the surface cleaning of specimens for \ac{EBSD} measurements, it is mostly known and used for materials other than cement. % and clay \cite{desbois.2017}.
For the preparation of cementitious specimens it is rarely applied.

\section{\acs*{ArBIB} polishing}
\label{sec:bib_pol_moving_result}

While the ion polishing is mostly used to enable \ac{EBSD}-measurements, it can also allow artefact free imaging of cementitious materials using the \ac{BSE} mode at low voltages (e.g., \qty{2}{\kilo\volt}).
This is demonstrated in \zcref{fig:BSE_CEMI_voltages} using a three year old hydrated CEM~I sample, which was first carbon coated and then cleaned using \ac{ArBIB} (\qty{6}{\kilo\volt} at \qty{2}{\milli\ampere} for \qty{20}{\minute}).
There is a significant difference in phase contrast, detail and susceptibility to charging artefacts at different voltages and with and without a carbon coating.
All images were recorded at \qty{16}{\bit} color depth.

\begin{figure*}[htb]
    \centering

	\includegraphics[width=\textwidth]{methods/bse-voltages/marked.pdf}
    \caption{
    	\ac{SEM} images of CEM~I hydrated for three years. With a carbon coating (top row) and after a subsequent \ac{ArBIB} polishing (middle row).Recorded using different voltages and at \qty{0.8}{\nano\ampere}.
		On the bottom, a comparison of the phase contrast between \ac{CH} and \ac{CSH} at \num{2.0} and \qty{10}{\kilo\volt} is illustrated.
    }
    \label{fig:BSE_CEMI_voltages}
\end{figure*}

There is no fixed grey value for a phase in \ac{BSE} imaging.
It can be manipulated in device not only by changing the voltage and current, but also by changing brightness, gamma and contrast values either automatically or manually.
Therefore, the contrast of the images was normalised, which in some instances significantly changed the appearance of the images.

The \ac{BSE} image of a carbon coated specimen mostly shows the coating surface at \qty{2.5}{\kilo\volt} and the chemical contrast is only slightly visible.
At \qty{4.0}{\kilo\volt} artefacts are still very prominent, making the image unusable for any segmentation task.
An artefact free image can be produced at \qty{10}{\kilo\volt}, allowing to distinguish between \ac{CH} and \ac{CSH} and between \ac{C3S} and \ac{C2S}.

An \ac{ArBIB} cleaned surface can be imaged at voltages as low as \qty{2.0}{\kilo\volt} (at \qty{0.8}{\nano\ampere}), resulting in an usable image.
The microstructure also appears more detailed than in images of the coated surface at \qty{10}{\kilo\volt}.
This is caused by the smaller interaction volume at lower voltages (see \zcref{fig:penetration_depth})
However, the phase contrast between \ac{CH} and \ac{CSH} is diminished compared to the carbon coated specimens at \qty{15}{\kilo\volt}.

The images of the cleaned surfaces suffer from charging artefacts starting at \qty{4.0}{\kilo\volt}, \qty{0.8}{\nano\ampere}.
At \qty{10}{\kilo\volt}, \qty{0.8}{\nano\ampere}, the image becomes unusable.
While current reduction to \qty{50}{\pico\ampere} eliminates the visible charging, the signal becomes very noisy.
Therefore, much longer frame integration times are required, which makes it difficult to obtain sharp details, since images easily become unstable and blurry due to the charging.

Furthermore, the \ac{ArBIB} cleaned surface shows more cracks due to dry-shrinking caused either by the heat from the ion beam or the prolonged vacuum exposure.

If a sample is not coated in the first place, \ac{ArBIB} polishing removes oil residue and reduces scratches as visible in \zcref{fig:mech_polishing_damage}.

\section{\acs*{ArBIB} sectioning}
\label{sec:bib_pol_fixed_results}

\ac{ArBIB} milling allows to image porous specimens without the requirement to embed them in epoxy resin (see \zcref{fig:BSE_ArBIB_CEMI}).
This is especially useful for specimens with pronounced capillary porosity.
Mechanical polishing would damage this not stabilized microstructure and the pores would be filled with debris.

\begin{figure*}[htb]
    \centering

	\includegraphics[width=\textwidth]{literature/ArBIB-CEMI-BSE.pdf}
    \caption{
       \ac{SEM} image of \qty{28}{\day} hydrated CEM~I.
       The larger left grain consists of multiple clinker phases, which can hardly be differentiated from each other.
       The surface was polished using \acl*{ArBIB} and was captured using a through the lens detector in \ac{BSE} mode at \qty{2.0}{\kilo\volt} and \qty{0.1}{\nano\ampere}.
       The contrast was increased in post processing.
    }
    \label{fig:BSE_ArBIB_CEMI}
\end{figure*}


The ion beam milled specimens have to be imaged using low voltages, since the so prepared surface is intrinsically non conductive and a consecutive coating would also alter the microstructure (see \zcref{fig:coating-influence}).
As demonstrated in the previous section, low voltages also result in a limited contrast between phases like \ac{CH} and \ac{CSH} or \ac{C3S} and \ac{C2S}.

Nevertheless, this method of preparation allows a close-to-native imaging of cementitous materials and results in a very flat and clean surface.

% TODO: hier auf die Nanoporosität eingehen, die freigelegt wird.

\section{\acs{ArBIB} preparation for nanoindentation combined with \ac{SEM} - \ac{BSE} and other techniques}

Nanoindentation is a technique which requires very flat surfaces.
While mechanical polishing can work if done with an oil-based polishing process, \ac{ArBIB} is a promising preparation method, since it does not require resin embedding, avoids chemical contamination and minimises mechanical damage as demonstrated in the section above.

As shown in \zcref{fig:lm_ni}, an usable area of \qty{7.5}{\milli\meter\squared} was prepared using \ac{ArBIB} sectioning.
The polishing was done at \qty{6}{\kilo\volt}, \qty{2}{\milli\ampere} in multiple sessions for a total of $\approx$ \qty{3}{\day}.
Between the sessions, the specimen had to be repositioned causing a damaged front area.

\begin{figure}[htb!]
  \centering
  \includegraphics[width=\linewidth]{nanoindentation/LM_overview.pdf}
  \caption{
    Light micrograph of a mortar prism (\qtyproduct{9x3.5}{\milli\meter}) which was sectioned using \ac{ArBIB} within \qty{72}{\hour}.
    The ion beams originated from the bottom, and a usable area of $\approx$ \qty{7.5}{\milli\meter\squared} was generated.
    On the top right, an image of the progress after \qty{48}{\hour} is shown during the process.
    The violet hue on the tungsten blade results from the ion beam.
  }
  \label{fig:lm_ni}
\end{figure}

The objective was to obtain a clean and smooth \ac{ITZ} between the binder and the quartz grain in the centre of the image.
To test the surface quality, a height map was measured using a confocal laser scanning microscope (\textsc{Keyence VK-X3050}).
The phase segmentation was enabled by an element mapping using \ac{mXRF}.
The quartz grains were segmented based on the \ce{Si} signal.
Furthermore, a \ac{BSE} image was acquired to find a suitable \ac{ROI} for the \ac{EDX} analysis and the nanoindentation sites.
\zcref{fig:arbib-experiment-positions} illustrates the position of the different areas measured by these methods.

\begin{figure}[htb!]
  \centering
  \includegraphics[width=\linewidth]{nanoindentation/measurement_overview.pdf}
  \caption{
    Illustration of the measurement positions of the different methods.
    The coloured area indicates the acquired height map (here aligned to the unpolished black surface) and the ROI which is referenced later (white dashed rectangle).
    {\color{cyan}Cyan}: \ac{BSE} image; {\color{red}red}: \ac{EDX} and nanoindentation area; thin black outlines: Quartz grains based on the \ac{mXRF} segmentation.
  }
  \label{fig:arbib-experiment-positions}
\end{figure}

The raw height map in \zcref{fig:arbib-experiment-positions} was aligned to the black, unpolished surface (compare \zcref{fig:lm_ni}) using three points.
This plane was defined as the zero level.
The aligned height map indicates that the ion beam milled the sample at an angle of \qtyrange{5}{6}{\degree} relative to the top surface.
Defining the polished surface as a reference plane using three points reveals the unevenness caused by the three beams (see \zcref{fig:arbib-topography}).
The central area was milled deeper than the lateral areas, creating a V-shaped pattern.
The polished area in the selected \ac{ROI} has a total height difference of approximately \qty{75}{\micro\meter}, whereas local variations are significantly smaller.

\begin{figure}[htb!]
  \centering

  \includegraphics[width=.8\linewidth]{nanoindentation/topography.pdf}
  \caption{\ac{ROI} of the height map, aligned to the \ac{ArBIB} polished surface.}
  \label{fig:arbib-topography}
\end{figure}

\textsc{ISO 25178} \cite{ISO25178-2:2022} defines roughness parameters like $S_{dr}$.
The developed area ratio $S_{dr}$ is the percentage increase in surface area due to texture compared to a perfectly flat surface.
It is \num{0} for perfectly smooth surfaces and increases with increasing roughness.
This parameter was calculated within a \qtyproduct{20 x 20}{px} sliding window for the whole \ac{ROI}.
The result is shown in \zcref{fig:arbib-rough}.
Combined with the segmentation results from the \ac{mXRF} dataset, this allows the surface roughness to be attributed to different phases.

\begin{figure}[htb!]
  \centering
  
  \includegraphics[width=.8\linewidth]{nanoindentation/local-roughness-Sdr.pdf}
  \caption{Developed area ratio $S_{dr}$ (right) with overlaid quartz aggregate positions determined by \ac{mXRF}.}
  \label{fig:arbib-rough}
\end{figure}

Calculating the mean developed area ratio for the aggregate area and the binder area results in \qty{59.7(33.5)}{\percent}, and \qty{87.8(41.5)}{\percent} respectively.
In other words, the the quartz grains have a smoother surfaces than the cement binder phase.
This is caused by the homogeneous composition of and the few phase boundaries (e.g. pores) within the quartz grains.
The mean developed area ratios calculated for the aggregate and binder areas were \qty{59.7(33.5)}{\percent} and \qty{87.8(41.5)}{\percent}, respectively.
In other words, the quartz grains exhibit a smoother surface than the cement binder phase.
This can be attributed to the homogeneous composition and the scarcity of phase boundaries (e.g., pores) within the quartz grains, while the binder matrix is full of small phase boundaries and incorporates lots of porosity.

A large scale \ac{SEM}-\ac{BSE} image was made to find suitable \ac{EDX} and nanoindentation areas and to estimate the porosity of the \ac{ITZ}.




Aligning these images from multiple modalities can be helpful to find a fitting measurement position but also to evaluate the results in a later stage.
The alignment can be done manually as discussed in the previous section by finding at least three or more matching point pairs or by manually transforming the images until the remaining error is acceptable.
However, a more automatic method would be preferred.
Multimodal image matching is a known but difficult problem in computer vision.
Widley used algorithms like \ac{SIFT} usually expect images of the same modality and often fail if the modality changes.
One possible solution is to bring all images in a similar modality.
For example by first applying denoising and edge detection algorithms.



This significantly helps to




nanoindentation: \cite{li.2025, li.2026}